\subsection{COMPARISON OF FILTERS}

DEFINITION 2. Given two filters \(\mathfrak{F}\) , \(\mathfrak{F}'\) on the same set X, \(\mathfrak{F}'\) is said to be finer than \(\mathfrak{F}\) , or \(\mathfrak{F}\) is coarser than \(\mathfrak{F}'\) , if \(\mathfrak{F} \subset \mathfrak{F}'\) . If also \(\mathfrak{F} \neq \mathfrak{F}'\) , then \(\mathfrak{F}'\) is said to be strictly finer than \(\mathfrak{F}\) , or \(\mathfrak{F}\) strictly coarser than \(\mathfrak{F}'\) .

Two filters are said to be comparable if one is finer than the other. The set of all filters on X is ordered by the relation ``F is coarser than F''; this relation is induced by the inclusion relation in \(\mathfrak{P}(\mathfrak{P}(\mathbf{X}))\) .

Let \((\mathfrak{F}_t)_{\iota \in I}\) be any non-empty family of filters on a set \(X\) (which must therefore be non-empty); then the set

\[
\mathfrak {F} = \bigcap_ {i \in I} \mathfrak {F} _ {i}
\]

satisfies axioms \((\mathbf{F}_{\mathbf{I}})\) , \((\mathbf{F}_{\mathbf{II}})\) and \((\mathbf{F}_{\mathbf{III}})\) and is therefore a filter; F is called the intersection of the family of filters \((\mathfrak{F}_{i})_{i\in\mathbf{I}}\) and is obviously the greatest lower bound of the set of the \(F_{i}\) in the ordered set of all filters on X.

The filter formed by the single set X is the smallest element of the ordered set of all filters on X. We shall see in no. 4 that, if X has more than one element, the set of all filters on X has no greatest element.

Given a set \(\mathfrak{G}\) of subsets of a set X, let us consider whether there are any filters on X which contain \(\mathfrak{G}\) . If such a filter exists then by \((\mathrm{F}_{\mathrm{II}})\) it contains also the set \(\mathfrak{G}'\) of finite intersections of sets of \(\mathfrak{G}\) (including X, which is the intersection of the empty subset of \(\mathfrak{G}\) ); hence a necessary condition for such a filter to exist is that the empty subset of X is not in \(\mathfrak{G}'\) . This condition is also sufficient, for by \((\mathrm{F}_{\mathrm{I}})\) any filter which contains \(\mathfrak{G}'\) also contains the set \(\mathfrak{G}''\) of subsets of X which contain a set of \(\mathfrak{G}'\) . Now \(\mathfrak{G}''\) clearly satisfies \((\mathrm{F}_{\mathrm{I}})\) ; it satisfies \((\mathrm{F}_{\mathrm{II}})\) by reason of the definition of \(\mathfrak{G}'\) ; and finally it satisfies \((\mathrm{F}_{\mathrm{III}})\) because the empty subset of X does not belong to \(\mathfrak{G}'\) . Hence \(\mathfrak{G}''\) is the coarsest filter which contains \(\mathfrak{G}\) , and we have proved:

PROPOSITION 1. A necessary and sufficient condition that there should exist a filter on X containing a set \(\mathfrak{G}\) of subsets of X is that no finite subset of \(\mathfrak{G}\) has an empty intersection.

The filter \(\mathfrak{G}''\) defined above is said to be generated by \(\mathfrak{G}\) , and \(\mathfrak{G}\) is said to be a subbase of \(\mathfrak{G}''\) .

Example. Let \(\mathfrak{S}\) be any set of subsets of a set X, and let \(\mathfrak{C}\) be the topology on X generated by \(\mathfrak{S}\) ( \(\S\) 2, no. 3, Example II). Since the set of finite intersections of sets of \(\mathfrak{S}\) is a base of \(\mathfrak{C}\) , it follows from the proof of Proposition 1 above and from Proposition 3 of \(\S\) 1, no. 3 that for each \(x \in X\) the neighbourhood filter of \(x\) for \(\mathcal{C}\) is generated by the set \(\mathfrak{S}(x)\) of sets of \(\mathfrak{S}\) which contain \(x\) .

COROLLARY 1. Let \(\mathfrak{F}\) be a filter on a set X, and A a subset of X. Then there is a filter \(\mathfrak{F}'\) which is finer than \(\mathfrak{F}\) and such that \(A \in \mathfrak{F}'\) , if and only if A meets all the sets of \(\mathfrak{F}\) .

COROLLARY 2. A set \(\Phi\) of filters on a non-empty set \(\mathbf{X}\) has a least upper bound in the set of all filters on \(\mathbf{X}\) if and only if, for all finite sequences \((\mathfrak{F}_i)_{1\leqslant i\leqslant n}\) of elements of \(\Phi\) and all \(\mathrm{A}_i\in \mathfrak{F}_i\) ( \(1\leqslant i\leqslant n\) ), the intersection \(\mathrm{A}_1\cap \dots \cap \mathrm{A}_n\) is not empty.

For this condition expresses that the union \(\mathfrak{G}\) of the filters \(\mathfrak{F} \in \Phi\) satisfies the condition of Proposition 1.

COROLLARY 3. The ordered set of all filters on a non-empty set X is inductive.

For every linearly ordered set \(\Phi\) of filters on \(X\) satisfies the condition of Corollary 2 of Proposition 1, since the sets \(A_{i}\) all belong to the same \(\mathfrak{F}_{j}\) by hypothesis, and we can apply \((\mathbf{F}_{\mathrm{II}})\) .

